\section{Results}
\label{sec:results}

This section details the empirical findings from our $N=108$ balanced simulation runs, computed via $10,000$ bootstrap resamples at a $95\%$ confidence interval (seed $= 2083521328$).
We evaluate the operational dynamics of our interactive paradigms against their respective uninstructed baselines (Stage 1) and conduct an absolute comparison between the question types (Stage 2).

\subsection{Stage 1: Interactive Conditions vs. Uninstructed Baselines}
The baseline conditions bypass conversational alignment entirely, meaning they register zero values for turns, dialogue tokens, and user response text. 
The primary value of the baseline lies in establishing a control threshold for downstream cluster geometry.

As summarized in Table~\ref{tab:results_summary}, every interactive question type required a statistically significant volume of conversational exchange compared to the instantaneous control baseline.
However, evaluating the secondary outcome reveals an unexpected geometric convergence.
Across all three tested formats—yes/no, multiple-choice, and open-ended—the resulting mean faithfulness scores do not manifest a statistically significant delta over their respective uninstructed control baselines. 

The bootstrapped $95\%$ confidence intervals for the mean difference $(\text{Condition} - \text{Baseline})$ all encompass zero:
\begin{itemize}
    \item \textbf{Yes/No Faithfulness Delta:} 
    $+0.222$, $95\%$ CI: $[-0.148, +0.602]$ (Not Significant).
    \item \textbf{Multiple-Choice Faithfulness Delta:}
    $+0.037$, $95\%$ CI: $[-0.333, +0.389]$ (Not Significant).
    \item \textbf{Open-Ended Faithfulness Delta:}
    $+0.231$, $95\%$ CI: $[-0.120, +0.593]$ (Not Significant).
\end{itemize}
This indicates that within the parameters of this pipeline, the geometric separation and alignment modifications generated by the instruction-tuned embedding layer did not yield a mathematically distinct advantage over blind baseline assignments. 

To provide a more granular diagnostic look into these overall performance patterns, Table~\ref{tab:submetrics_breakdown} isolates and breaks down the three underlying dimensions—coherence, alignment, and separation—that compile the aggregate mean faithfulness ratings.

% =========================================================================
%  TABLE 1: PRIMARY DIALOGUE DYNAMICS & MEAN FAITHFULNESS
% =========================================================================
\newsavebox{\tableonebox}

\begin{table*}[t]
\centering
\sbox{\tableonebox}{%
\begin{tabular}{lcccc}
\toprule
\textbf{\makecell[l]{Condition \\ Framework}} & 
\textbf{\makecell[c]{Turns to \\ Convergence}} & 
\textbf{\makecell[c]{Total \\ Tokens}} & 
\textbf{\makecell[c]{User \\ Tokens}} & 
\textbf{\makecell[c]{Mean \\ Score}} \\
\midrule
\textbf{Yes/No Condition}    & $12.97$ & $1238.81$ & $338.69$ & $2.20$ \\
\textit{Yes/No Baseline}     & \textit{0.00} & \textit{0.00} & \textit{0.00} & \textit{1.98} \\
\midrule
\textbf{Multiple-Choice}     & $\mathbf{7.83}$ & $\mathbf{1196.42}$ & $\mathbf{100.92}$ & $\mathbf{2.23}$ \\
\textit{Multiple-Choice Base}& \textit{0.00} & \textit{0.00} & \textit{0.00} & \textit{2.19} \\
\midrule
\textbf{Open-Ended Condition} & $9.64$ & $1383.81$ & $616.50$ & $2.19$ \\
\textit{Open-Ended Baseline}  & \textit{0.00} & \textit{0.00} & \textit{0.00} & \textit{1.96} \\
\bottomrule
\end{tabular}%
}
\begin{minipage}{\wd\tableonebox}
\centering
\usebox{\tableonebox}
\vspace{0.2cm}
\caption{Aggregated Dialogue Dynamics and Mean Faithfulness Scores Across Interactive Frameworks ($N=108$ Runs)}
\label{tab:results_summary}
\end{minipage}
\end{table*}


% =========================================================================
%  TABLE 2: GRANULAR SUB-METRICS BREAKDOWN
% =========================================================================
\newsavebox{\tabletwobox}

\begin{table*}[t]
\centering
% Slightly increase internal cell padding for Table 2 so its 4 columns 
% naturally match the overall visual presence of Table 1 beautifully
\setlength{\tabcolsep}{8pt} 
\sbox{\tabletwobox}{%
\begin{tabular}{lccc}
\toprule
\textbf{\makecell[l]{Condition \\ Framework}} & 
\textbf{\makecell[c]{Coherence \\ Score}} & 
\textbf{\makecell[c]{Alignment \\ Score}} & 
\textbf{\makecell[c]{Separation \\ Score}} \\
\midrule
\textbf{Yes/No Condition}    & $2.58$ & $1.75$ & $2.28$ \\
\textit{Yes/No Baseline}     & \textit{2.44} & \textit{1.58} & \textit{1.92} \\
\midrule
\textbf{Multiple-Choice}     & $\mathbf{2.67}$ & $1.72$ & $2.31$ \\
\textit{Multiple-Choice Base}& \textit{2.47} & \textit{\textbf{1.78}} & \textit{\textbf{2.33}} \\
\midrule
\textbf{Open-Ended Condition} & $2.61$ & $1.75$ & $2.22$ \\
\textit{Open-Ended Baseline}  & \textit{2.44} & \textit{1.42} & \textit{2.03} \\
\bottomrule
\end{tabular}%
}
\begin{minipage}{\wd\tabletwobox}
\centering
\usebox{\tabletwobox}
\vspace{0.2cm}
\caption{Granular Breakdown of LLM-as-Judge Faithfulness Sub-Metrics}
\label{tab:submetrics_breakdown}
\end{minipage}
\end{table*}

\subsection{Stage 2: Pairwise Question Type Comparison}
To determine the optimal human-computer interface, we analyze the absolute variations between the interactive treatments across our observed criteria.

\subsubsection{Turns to Convergence (Primary Outcome)}
The primary efficiency metric reveals a highly distinct separation between paradigms, isolating an absolute hierarchical rank. 
Multiple-choice questions proved to be the most efficient strategy, converging in a mean of just $7.83$ turns. 
Open-ended queries followed with a mean of $9.64$ turns, while yes/no interactions required the most extensive dialogues, averaging $12.97$ turns to convergence.

Pairwise bootstrapping confirms the statistical significance of this ranking:
\begin{itemize}
    \item \textbf{Yes/No vs. Multiple-Choice:} 
    Multiple-choice dramatically outperforms yes/no ($\text{Diff} = +5.139$, $95\%$ CI: $[+3.111, +7.250]$). 
    \item \textbf{Yes/No vs. Open-Ended:} 
    Open-ended queries converge significantly faster than binary choices ($\text{Diff} = +3.333$, $95\%$ CI: $[+1.028, +5.583]$).
    \item \textbf{Multiple-Choice vs. Open-Ended:} 
    Multiple-choice establishes a statistically verified lead over open-ended configurations ($\text{Diff} = -1.806$, $95\%$ CI: $[-3.556, -0.139]$).
\end{itemize}

\subsubsection{Token Consumption and Question Clarity}
The volume of language exchanged tracks the structural density of the interaction profiles. 
Total tokens exchanged exhibited no statistically significant variation between the yes/no framework ($1238.81$) and the multiple-choice framework ($1196.42$), with a bootstrapped difference interval of $[-241.974, +337.619]$, which strictly encompasses zero. 
Similarly, open-ended token generation ($1383.81$) fell within a non-significant baseline margin compared to its peers.

However, the user response token count—serving as our direct proxy for cognitive load and question clarity—reveals massive, statistically significant variations. 
Multiple-choice questions required a sparse mean of $100.92$ tokens from the user.
Yes/no questions demanded significantly more descriptive clarifications from the simulated agent, averaging $338.69$ tokens.
Open-ended prompts generated the heaviest linguistic overhead, forcing the user to draft an extensive mean of $616.50$ tokens per run. 

Pairwise testing confirms that multiple-choice queries minimize user token generation over yes/no ($\text{Diff} = +237.778$, $95\%$ CI: $[+150.722, +332.640]$), while yes/no interfaces restrict response lengths significantly more effectively than unstructured open questions ($\text{Diff} = -277.806$, $95\%$ CI: $[-406.444, -151.358]$).

\subsubsection{Downstream Faithfulness Alignment}
In contrast to the highly dynamic efficiency metrics, the resulting mean faithfulness scores remain entirely uniform across the interactive strategies. 
The composite mean scores reached $2.20$ for yes/no, $2.23$ for multiple-choice, and $2.19$ for open questions.
The corresponding sub-metrics—coherence, alignment, and separation—maintain flat distributions across all parameters. 
Pairwise bootstrapping verifies that changing the question style yields absolutely no statistical difference in final cluster faithfulness.

\subsection{Evaluation Framework Validation: Inter-Rater Agreement}
\label{subsec:inter_rater_agreement}

To assess the reliability and operational validity of our downstream LLM-as-a-judge framework, a human-in-the-loop cross-validation check was executed.
A human expert independently annotated a random sample of clustering configurations ($N = 6$ items, yielding $18$ total individual metric ratings) across the three core evaluation dimensions. 
Inter-rater agreement between the human expert and the automated LLM judge was quantified using a quadratically weighted Cohen's Kappa ($\kappa$).
To mathematically account for the restricted sample volume, a $10,000$-resample bootstrap was applied to calculate 95\% confidence intervals.
The agreement metrics are detailed in Table~\ref{tab:cohens_kappa}.

% =========================================================================
%  TABLE 3: INTER-RATER AGREEMENT (COHEN'S KAPPA)
% =========================================================================
\newsavebox{\tablethreebox}

\begin{table*}[t]
\centering
\sbox{\tablethreebox}{%
\begin{tabular}{lccc}
\toprule
\textbf{Evaluation Metric} & \textbf{Items Rated} & \textbf{Weighted Cohen's $\kappa$} & \textbf{95\% Bootstrap CI} \\ 
\midrule
~~Mean Score\textsuperscript{a} & 6 & 0.571 & $[-0.667, +0.909]$ \\ 
~~Coherence & 6 & 0.690 & $[+0.000, +0.889]$ \\
~~Alignment & 6 & 0.640 & $[+0.000, +0.909]$ \\
~~Separation & 6 & 0.385 & $[-0.500, +0.769]$ \\ 
\bottomrule
\multicolumn{4}{l}{\footnotesize \textsuperscript{a} Calculated as the composite mean: $(\text{coherence} + \text{alignment} + \text{separation}) / 3$}
\end{tabular}%
}
\begin{minipage}{\wd\tablethreebox}
\centering
\setcounter{table}{2} % Explicitly forces this to be numbered as Table 3
\usebox{\tablethreebox}
\vspace{0.2cm}
\caption{Inter-Rater Agreement Metrics (Weighted Cohen's $\kappa$) between Human Expert and LLM Judge}
\label{tab:cohens_kappa}
\end{minipage}
\end{table*}

The primary composite Mean Score yielded a $\kappa = 0.571$, indicating a \textit{moderate} level of baseline agreement between human judgment and the automated evaluator. 
Evaluated individually, the system exhibited \textit{substantial} agreement on foundational text traits such as semantic coherence ($\kappa = 0.690$) and profile alignment ($\kappa = 0.640$).
Conversely, spatial separation achieved only \textit{fair} agreement ($\kappa = 0.385$), capturing the inherent subjectivity involved in contextually isolating overlapping technical terms. 

Due to the high operational cost of expert manual annotations, this validation was strictly restricted to an exploratory sample ($N = 6$). 
Consequently, the resulting bootstrap confidence intervals are expectedly wide, and these coefficients must be interpreted strictly as an indicative sanity check of judge behavior rather than an absolute population statistic.
Nevertheless, the moderate agreement observed across the primary score configurations confirms that the automated evaluator operates as a functional proxy for tracking relative performance trends across our primary simulation matrix.